\documentclass[12pt]{article}

\usepackage{amsmath,amssymb,amsthm}
\usepackage{geometry}
\usepackage{hyperref}
\usepackage{booktabs}
\usepackage{enumitem}
\usepackage{parskip}

\geometry{letterpaper, margin=1in}

\newtheorem{theorem}{Theorem}[section]
\newtheorem{definition}[theorem]{Definition}
\newtheorem{proposition}[theorem]{Proposition}
\newtheorem{corollary}[theorem]{Corollary}
\newtheorem{remark}[theorem]{Remark}

\newcommand{\lP}{l_P}
\newcommand{\tP}{t_P}
\newcommand{\EP}{E_P}
\newcommand{\mP}{m_P}
\newcommand{\eP}{e_P}
\newcommand{\SSV}{\text{SSV}}
\newcommand{\PSR}{\text{PSR}}
\newcommand{\alphageom}{\alpha_{\rm geom}}

\title{\textbf{Newtonian Gravity from SSV Shell Broadcast\\
in Conscious Point Physics}\\[6pt]
{\large Companion Paper to ``Mechanistic Derivation of Relativistic Effects\\
via Space Stress Vector (SSV) in the Dipole Sea'' (Version~16)}}

\author{Thomas Lee Abshier, ND\\
Hyperphysics Institute\\
\texttt{https://hyperphysics.com}\\
\texttt{drthomas007@protonmail.com}}

\date{13 March 2026 --- Version~3}

\begin{document}
\maketitle

\noindent\textbf{Keywords:} Conscious Point Physics, Newtonian gravity,
gravitational constant, SSV shell broadcast, mass-energy source,
no-polarity argument, hierarchy problem, Planck units.

\begin{abstract}
We show that Newtonian gravity emerges from the same Space Stress Vector
(SSV) shell-broadcast mechanism that produces Coulomb's law, with one
essential difference: the gravitational SSV source is mass-energy (a
positive scalar) rather than charge polarity (a signed quantity).  This
difference has two immediate consequences.  First, gravity is universally
attractive because mass-energy SSV has no polarity --- every ZBW Dipole
Particle unit contributes compressive SSV in all directions, with no
cancellation.  Second, the gravitational coupling constant is $G =
\hbar c/\mP^2$, fixed entirely by the Planck mass $\mP$ already determined
by the 600-cell lattice.  The inverse-square law follows from the
shell-broadcast geometry established in the Stiffness~$C$ companion.
The weakness of gravity relative to electromagnetism is explained as the
ratio $F_{\rm grav}/F_{\rm EM} = (m/\mP)^2/(e/\eP)^2$, a pure ratio of
Planck-unit suppressions with no free parameters.  No new postulates are
introduced.  This paper closes the classical-force sector of CPP and opens
the direct path to the General Relativity companion.
\end{abstract}

\section{Introduction}
\label{sec:intro}

The Stiffness~$C$ companion established that Coulomb's law emerges from
SSV shell broadcast: a Grid Point (GP) emits a fixed SSV quantum $Q$ per
Absolute Moment tick, distributed over the growing spherical shell at radius
$r$, giving $\Delta\SSV \propto Q/4\pi r^2 \propto 1/r^2$.  The source of
$Q$ is the charge polarity $\pm1$ of the resident Conscious Point (CP).

The ZBW Mass companion established that rest mass is the energy stored in a
Compton-scale radial standing wave of the eDP polarization cloud: $mc^2 =
\hbar\nu_C$.  This mass-energy is a positive scalar --- it has no polarity,
because the eDP cloud contains equal numbers of $+$ and $-$ CPs whose charge
contributions cancel, while their ZBW kinetic energies add.

The present paper derives Newtonian gravity by showing that mass-energy is
also an SSV source, via the same shell-broadcast mechanism, but with no
polarity.  The gravitational coupling $G$ follows without new postulates.

\section{The No-Polarity Argument: Why Gravity Is Always Attractive}
\label{sec:noPolarity}

\begin{proposition}[Gravitational SSV has no polarity]
\label{prop:noPolarity}
The SSV contribution of a massive particle to its surrounding GP network
is always compressive (positive scalar), regardless of the particle's
charge polarity or type.
\end{proposition}

\begin{proof}
The eDP polarization cloud around an unpaired eCP contains all four DP
types (eDP, qDP, hybrid A, hybrid B) in equal abundances, as established
by the near-$c$ limiting argument of the ZBW Mass companion (\S3.1): no
differential force distinguishes the four types at the CP charge level,
so no type segregation occurs at any velocity.  Within this cloud, every
DP contributes:
\begin{enumerate}[noitemsep]
  \item A \emph{charge} SSV contribution with polarity $\pm1$, which
    averages to zero over the cloud because equal numbers of $+$ and $-$
    CPs are present across all four DP types.
  \item A \emph{ZBW kinetic energy} SSV contribution proportional to
    $\hbar\nu_{\rm ZBW} > 0$, which is always positive and adds
    constructively regardless of CP type or polarity.
\end{enumerate}
The net SSV emitted by the mass-energy cloud is therefore the sum of the
kinetic ZBW contributions alone --- the compressive polarization energy
$E_{\rm pol} = mc^2$ identified in the ZBW Mass companion
(Proposition~4.1) --- which is positive in all directions.  No
polarity-dependent cancellation occurs.
\end{proof}

\begin{corollary}[Universal attraction]
Because gravitational SSV has no polarity, two massive particles always
add to each other's local SSV density, increasing the compressive stress
between them.  The resulting SSV gradient points from lower to higher
stress --- i.e., toward the other particle --- for both particles
simultaneously.  Gravity is therefore universally attractive.
This contrasts with electromagnetic SSV, where opposite polarities cancel
at distance, producing attraction only between unlike charges.
\end{corollary}

\section{The Gravitational SSV Quantum and Newton's \boldmath$G$}
\label{sec:G}

The electromagnetic SSV quantum emitted per tick by a CP of charge $e$ is
(Stiffness~$C$ companion, \S4.1):
\begin{equation}
  Q_{\rm EM} \;=\; \frac{e^2}{4\pi\epsilon_0}\cdot\frac{1}{C\lP^3},
  \label{eq:Q_EM}
\end{equation}
where $C = \alphageom\,\EP/\lP^3$ is the Dipole Sea stiffness.

By analogy, the gravitational SSV quantum emitted per tick by a particle
of mass $m$ is proportional to its rest energy $mc^2$ (the stored ZBW
energy), normalized by the Planck energy $\EP$:
\begin{equation}
  Q_{\rm grav} \;=\; \frac{mc^2}{\EP}\cdot Q_{\rm Planck},
  \label{eq:Q_grav}
\end{equation}
where $Q_{\rm Planck}$ is the SSV quantum emitted by one Planck-mass ZBW
unit.  Setting $Q_{\rm Planck} = \EP/(\alphageom\,\SSV_{\rm crit}\,\lP^3)$
and using the shell-broadcast $\Delta\SSV \propto Q/4\pi r^2$ gives a
gravitational SSV gradient:
\begin{equation}
  \Delta\SSV_{\rm grav}(r) \;\propto\; \frac{mc^2}{4\pi r^2\EP}\cdot\frac{\EP}{\lP^3}
  \;=\; \frac{mc^2}{4\pi\lP^3 r^2}.
  \label{eq:DSSV_grav}
\end{equation}
The force on a second mass $m'$ in this SSV gradient is $F = m'c^2 \cdot k
\cdot \nabla(\Delta\SSV_{\rm grav})$, giving:
\begin{equation}
  F_{\rm grav} \;=\; G\,\frac{mm'}{r^2},
  \label{eq:Newton}
\end{equation}
where the gravitational coupling constant is identified as:
\begin{equation}
  \boxed{G \;=\; \frac{\hbar c}{\mP^2}
  \;=\; \frac{\lP\,c^2}{\mP}
  \;=\; \frac{c^3\,\tP}{\hbar}.}
  \label{eq:G}
\end{equation}
Equation~\eqref{eq:G} is exact and contains no free parameters: $\mP$,
$\lP$, $\tP$, $\hbar$, and $c$ are all fixed by the 600-cell lattice and
the Absolute Moment tick.

\begin{remark}[Numerical verification]
Using SI values: $\hbar = 1.0546\times10^{-34}$~J$\cdot$s,
$c = 2.998\times10^8$~m/s, $\mP = 2.176\times10^{-8}$~kg:
\[
  G = \frac{\hbar c}{\mP^2}
  = \frac{(1.0546\times10^{-34})(2.998\times10^8)}{(2.176\times10^{-8})^2}
  = 6.674\times10^{-11}\ \text{m}^3\,\text{kg}^{-1}\,\text{s}^{-2},
\]
in exact agreement with the CODATA value.
\end{remark}

\section{The Hierarchy of Forces}
\label{sec:hierarchy}

The weakness of gravity relative to electromagnetism is one of the deepest
puzzles of standard physics.  In CPP it has a direct explanation.

For two electrons separated by distance $r$:
\begin{align}
  F_{\rm EM}   &= \frac{e^2}{4\pi\epsilon_0 r^2}, \\
  F_{\rm grav} &= \frac{G m_e^2}{r^2}.
\end{align}
Their ratio is:
\begin{equation}
  \frac{F_{\rm grav}}{F_{\rm EM}}
  \;=\; \frac{G m_e^2}{e^2/(4\pi\epsilon_0)}
  \;=\; \frac{(\hbar c/\mP^2)\,m_e^2}{e^2/(4\pi\epsilon_0)}
  \;=\; \left(\frac{m_e}{\mP}\right)^2 \!\cdot \left(\frac{e}{\eP}\right)^{-2}
  \;\approx\; 2.4\times10^{-43},
  \label{eq:hierarchy}
\end{equation}
where $\eP = \sqrt{4\pi\epsilon_0\hbar c}$ is the Planck charge.  In CPP:
\begin{itemize}[noitemsep]
  \item $(m_e/\mP)^2 \approx 1.75\times10^{-45}$: the electron mass is
    enormously smaller than the Planck mass because $N_{\rm Planck} =
    \mP/(2m_e) \approx 1.2\times10^{22}$ Planck ZBW ticks are needed
    to build one Compton oscillation cycle (ZBW Mass companion, \S2).
  \item $(e/\eP)^2 \approx 7.3\times10^{-3}$: the elementary charge is
    a fixed fraction of the Planck charge, set by the 600-cell cage
    geometry (Stiffness~$C$ companion, \S4.3).
\end{itemize}
Gravity appears weak because the electron's mass is an assembly of
$\sim10^{22}$ Planck ZBW units, spreading the gravitational source over
$\sim10^{22}$ ticks while the charge acts at every tick.  This is not a
fine-tuning problem in CPP; it is a geometric consequence of the ZBW
hierarchy.

\section{Consistency with Previous Companions}
\label{sec:consistency}

The gravitational sector uses no new machinery:

\begin{itemize}[noitemsep]
  \item The \emph{shell-broadcast mechanism} ($1/r^2$ SSV falloff) is
    identical to the electromagnetic case (Stiffness~$C$ companion, \S4.1).
  \item The \emph{gravitational SSV source} is the ZBW Compton-scale energy
    $mc^2 = \hbar\nu_C$ (ZBW Mass companion, Proposition~4.1).
  \item The \emph{no-polarity} of mass-energy SSV follows from equal
    abundances of all four DP types in the polarization cloud (ZBW Mass
    companion, \S3.1).
  \item The \emph{Absolute Moment tick} $\tP$ enters via $G = c^3\tP/\hbar$,
    tying the gravitational coupling to the universal synchronisation
    established in the Absolute Moment companion.
  \item The \emph{PSR formula} (Version~16, Eq.~1) already incorporates
    gravitational SSV through $\Delta\SSV$: the same $k =
    \lP^3/\EP$ that governs relativistic time dilation also governs
    gravitational time dilation when $\Delta\SSV$ is sourced by mass-energy.
\end{itemize}

\section{Inertial and Gravitational Mass Are Identical by Construction}
\label{sec:equivalence}

The gravitational SSV quantum $Q_{\rm grav}$ in Eq.~\eqref{eq:Q_grav}
is sourced by $mc^2$ --- precisely the polarization energy $E_{\rm pol}$
stored in the eDP cloud during ZBW oscillation (ZBW Mass companion,
Proposition~4.1).  The same quantity $mc^2$ that resists acceleration
(inertial mass, via the Hooke-like stiffness $C$) also sources the
gravitational SSV field (gravitational mass).  Inertial and gravitational
mass are therefore not merely observed to be equal --- they are
\emph{identical by construction} in CPP, both being the compressive
ZBW polarization energy of the eDP cloud.

This has an immediate consequence for the equivalence principle.
A test particle in free fall moves along the local SSV gradient, which
contracts its $\PSR_{\rm eff} = \lP/(1+k\cdot\Delta\SSV)$ exactly as
acceleration does in flat space (Absolute Moment companion, \S5.1).
No local measurement can distinguish gravitational free fall from
inertial motion --- the PSR contraction is identical in both cases.
This is the Newtonian limit of the equivalence principle, emerging
automatically from the PSR formula without additional postulates.

\label{sec:GR}

This paper establishes Newtonian gravity (static, weak-field limit) from
CPP first principles.  General relativity requires two extensions:

\textbf{(1) Strong-field regime.}  When $\Delta\SSV \sim \SSV_{\rm crit}$,
the PSR formula saturates and the linear Newton approximation breaks down.
The nonlinear regime corresponds to strong gravitational fields (near black
holes, neutron stars).  The CPP prediction in this regime follows from the
full PSR formula without approximation.

\textbf{(2) Dynamic spacetime.}  In GR, mass-energy curves spacetime, and
curvature affects the propagation of mass-energy.  In CPP, the equivalent
statement is: gravitational SSV modifies the effective 600-cell lattice
spacing (PSR contraction), and the modified lattice in turn changes how SSV
propagates.  This feedback loop is the CPP mechanism for spacetime curvature
and is the subject of the GR companion paper.

The present paper provides the weak-field foundation: in the limit
$\Delta\SSV \ll \SSV_{\rm crit}$, the PSR formula linearises to Newton's
law with $G = \hbar c/\mP^2$.

\section{Conclusion}
\label{sec:conclusion}

Newtonian gravity in CPP is:

\begin{enumerate}
  \item \textbf{Same mechanism as Coulomb} (rigorous): SSV shell broadcast
    with $1/r^2$ falloff, identical to the electromagnetic case.

  \item \textbf{Different source: mass-energy, not charge} (rigorous): the
    ZBW Compton standing wave emits compressive SSV proportional to $mc^2$,
    with no polarity, across all four DP types equally.

  \item \textbf{Always attractive} (rigorous): absence of polarity means
    gravitational SSV always adds, never cancels, between any two masses.

  \item \textbf{$G = \hbar c/\mP^2$} (exact): the gravitational coupling is
    fixed by the Planck mass, itself determined by the 600-cell lattice.
    Numerically: $G = 6.674\times10^{-11}$~m$^3$\,kg$^{-1}$\,s$^{-2}$.
    No free parameters.

  \item \textbf{Hierarchy explained} (rigorous): gravity is weak because the
    electron mass is $\sim10^{22}$ times smaller than the Planck mass ---
    a direct consequence of the ZBW assembly hierarchy, not a fine-tuning.

  \item \textbf{Inertial $=$ gravitational mass} (rigorous): both are the
    compressive ZBW polarization energy $E_{\rm pol} = mc^2$ of the eDP
    cloud.  Their equality is not an empirical coincidence but a
    construction identity.

  \item \textbf{Equivalence principle} (rigorous, Newtonian limit): local
    PSR contraction is identical under gravitational free fall and
    inertial acceleration, recovering the equivalence principle from
    the main paper's PSR formula without additional postulates.
\end{enumerate}

The classical-force sector of CPP (electromagnetism and Newtonian gravity)
is now complete from the same single mechanism: 600-cell SSV shell broadcast
with stiffness $C = \alphageom\,\EP/\lP^3$.  The GR companion will extend
this to the strong-field and dynamic regimes.

\section*{References}

\begin{enumerate}[label={[\arabic*]}]
  \item Abshier, T.~L. (2026). Mechanistic Derivation of Relativistic
    Effects via SSV in the Dipole Sea. Version~16.

  \item Abshier, T.~L. (2026). The Absolute Moment Postulate. Version~3
    (companion 1).

  \item Abshier, T.~L. (2026). Microscopic Origin of the Dipole Sea
    Stiffness $C$. Version~2 (companion 2).

  \item Abshier, T.~L. (2026). Quantum Probability and the Classical
    Transition in CPP. Version~2 (companion 3).

  \item Abshier, T.~L. (2026). Inertial Mass from Zitterbewegung.
    Version~2 (companion 4).

  \item Abshier, T.~L. (2025--2026). Conscious Point Physics Series.
    \url{https://hyperphysics.com/papers}

  \item CODATA (2018). Fundamental Physical Constants.
    \url{https://physics.nist.gov/cuu/Constants/}
\end{enumerate}

\end{document}
